\documentclass[11pt]{article}

\usepackage[margin=1in]{geometry}
\usepackage{amsmath,amssymb,booktabs,array,hyperref}
\usepackage{xcolor}

\title{Success Rates and Literature Context for Crossing Sequence Memory in Kuramoto Networks}
\author{}
\date{}

\newcommand{\ii}{\mathrm{i}}
\newcommand{\conj}[1]{\overline{#1}}
\newcommand{\Tcal}{\mathcal{T}}


\usepackage[
    backend=biber,
    style=numeric,
    sorting=none
]{biblatex}

\addbibresource{literature.bib}


\begin{document}

\maketitle

\section{Benchmark Problem and Kuramoto Setup}

We study crossing sequence memory in a Kuramoto network with $N=64$
oscillators.  The oscillator state and stored phase patterns are
\begin{equation}
    z_i=e^{\ii\theta_i},
    \qquad
    \xi_i^\mu=e^{\ii\phi_i^\mu},
    \qquad
    i=1,\ldots,N .
\end{equation}
The complex overlap between the network state and stored pattern $\mu$ is
\begin{equation}
    m_\mu(z)
    =
    \frac{1}{N}\sum_{j=1}^N \conj{\xi_j^\mu}z_j .
\end{equation}

The autonomous phase dynamics are
\begin{equation}
    \dot{\theta}_i
    =
    g\,\Imag\!\left[h_i(z)\conj{z_i}\right],
    \qquad
    \omega_i=0 ,
\end{equation}
where $h_i(z)$ is the static fourth-order learned field and $g$ is a global
coupling gain.  Thus all sequential motion must be generated by the learned
coupling tensor itself.  In the autonomous benchmark runs there is no external
drive after initialization.  Each trial starts from a small perturbation of the
first pattern of the target sequence,
\begin{equation}
    \theta_i(0)
    =
    \arg \xi_i^{(s,1)}
    +
    \eta_i,
    \qquad
    \eta_i\sim\mathcal{N}(0,\sigma_{\mathrm{init}}^2).
\end{equation}

The benchmark sequences are generated as continuous random corruption paths,
not as unrelated random patterns.  For a prescribed corruption ratio $\rho$,
consecutive patterns satisfy approximately
\begin{equation}
    d(\xi_t,\xi_{t+1}) \approx \rho ,
\end{equation}
where for phase patterns we use the overlap distance
\begin{equation}
    d(\xi,\zeta)
    =
    1-
    \left|
    \frac{1}{N}
    \sum_{i=1}^N
    \conj{\xi_i}\zeta_i
    \right|.
\end{equation}
This is the phase-pattern analogue of a normalized mismatch distance.  A
crossing graph is constructed by generating the shared pattern or shared segment
once and reusing it exactly.  The incoming and outgoing branches are independent
corruption paths away from the shared node or segment.

We use three benchmark topologies:
\begin{enumerate}
    \item[(a)] \textbf{Two-branch crossing.}
    Two length-five sequences share one central pattern:
    \begin{equation}
        \xi^{(s,1)}
        \to
        \xi^{(s,2)}
        \to
        \xi^{\star}
        \to
        \xi^{(s,4)}
        \to
        \xi^{(s,5)},
        \qquad s=1,2 .
    \end{equation}

    \item[(b)] \textbf{Three-branch crossing.}
    Three length-five sequences share one central pattern:
    \begin{equation}
        \xi^{(s,1)}
        \to
        \xi^{(s,2)}
        \to
        \xi^{\star}
        \to
        \xi^{(s,4)}
        \to
        \xi^{(s,5)},
        \qquad s=1,2,3 .
    \end{equation}

    \item[(c)] \textbf{Two-pattern shared segment.}
    Two length-six sequences share the segment
    $\xi^{\star,1}\to\xi^{\star,2}$:
    \begin{equation}
        \xi^{(s,1)}
        \to
        \xi^{(s,2)}
        \to
        \xi^{\star,1}
        \to
        \xi^{\star,2}
        \to
        \xi^{(s,5)}
        \to
        \xi^{(s,6)},
        \qquad s=1,2 .
    \end{equation}
\end{enumerate}
\section{Learning Rules}

For a phase state $z_i=e^{\mathrm{i}\theta_i}$ and stored phase patterns
$\xi^\alpha$, we measure the complex overlap with pattern $\alpha$ by
\begin{equation}
m_\alpha(z)
=
\frac{1}{N}
\sum_{j=1}^N
\overline{\xi_j^\alpha}z_j .
\end{equation}
All rules below are autonomous static fourth-order fields of bidegree $(2,1)$:
they use two overlap factors and one conjugate overlap factor, and no external
drive or hidden state.

\paragraph{Self rule / autoassociative stabilizer $(\mu)$.}
This term stabilizes each stored pattern but does not by itself create a
sequence:
\begin{equation}
h_i^{\mathrm{self}}(z)
=
\sum_\mu
\xi_i^{\mu}
m_{\mu}(z)
\left|m_{\mu}(z)\right|^2 .
\end{equation}

\paragraph{Forward quartic rule / Markov-1 rule $(\mu\rightarrow\mu+1)$.}
This is the simplest sequence rule.  It pushes the state from pattern $\mu$ to
its immediate successor:
\begin{equation}
h_i^{+1}(z)
=
\sum_\mu
\xi_i^{\mu+1}
m_{\mu}(z)
\left|m_{\mu}(z)\right|^2 .
\end{equation}

\paragraph{Skip-two quartic rule $(\mu\rightarrow\mu+2)$.}
This rule pushes from pattern $\mu$ directly toward the second successor:
\begin{equation}
h_i^{+2}(z)
=
\sum_\mu
\xi_i^{\mu+2}
m_{\mu}(z)
\left|m_{\mu}(z)\right|^2 .
\end{equation}

\paragraph{Mixture rules.}
We also test linear mixtures of the self, forward, and skip-two rules:
\begin{equation}
h_i^{\mathrm{mix}}(z)
=
a_0 h_i^{\mathrm{self}}(z)
+
a_1 h_i^{+1}(z)
+
a_2 h_i^{+2}(z).
\end{equation}
The tested cases include $\mu+\mu+1$, $\mu+\mu+2$,
$\mu+\mu+1+\mu+2$, and convex mixtures of $\mu+1$ and $\mu+2$.

\paragraph{Magnitude-gated bigram rule / one-step context rule.}
This is the simplest context rule.  The next state is driven by the current
overlap $m_{\mu+1}$, but gated by the residual overlap with the previous pattern
$m_\mu$:
\begin{equation}
h_i^{\mathrm{bigram}}(z)
=
\sum_\mu
\xi_i^{\mu+2}
m_{\mu+1}(z)
\left|m_{\mu}(z)\right|^2 .
\end{equation}
Equivalently, after reindexing,
\begin{equation}
h_i^{\mathrm{bigram}}(z)
=
\sum_\mu
\xi_i^{\mu+1}
m_{\mu}(z)
\left|m_{\mu-1}(z)\right|^2 .
\end{equation}

\paragraph{Multilag context rule.}
This extends the bigram rule by allowing several recent overlaps to gate the
transition:
\begin{equation}
h_i^{\mathrm{multilag}}(z)
=
\sum_\mu
\xi_i^{\mu+1}
m_\mu(z)
\left(
w_0 |m_\mu(z)|^2
+
w_1 |m_{\mu-1}(z)|^2
+
w_2 |m_{\mu-2}(z)|^2
\right).
\end{equation}

\paragraph{Coherent bigram rule.}
This is a phase-sensitive version of the bigram rule.  It uses the relative
phase between the current and previous overlaps, not only the magnitude of the
previous overlap:
\begin{equation}
h_i^{\mathrm{coh2}}(z)
=
\sum_\mu
\xi_i^{\mu+1}
m_\mu(z)^2
\overline{m_{\mu-1}(z)} .
\end{equation}

\paragraph{Coherent trigram rule.}
This is the minimal interpretable rule that uses two-step context while staying
within the same fourth-order tensor budget:
\begin{equation}
h_i^{\mathrm{coh3}}(z)
=
\sum_\mu
\xi_i^{\mu+1}
m_\mu(z)
m_{\mu-2}(z)
\overline{m_{\mu-1}(z)} .
\end{equation}
This rule can distinguish a shared segment because it can still read which
branch the trajectory came from two steps earlier.

\paragraph{Residual bigram rule.}
For a branching state $\mu$, let $\mathcal{S}(\mu)$ be the set of possible
successors and define the mean successor pattern
\begin{equation}
\bar{\xi}^{\mu+1}
=
\frac{1}{|\mathcal{S}(\mu)|}
\sum_{\nu\in\mathcal{S}(\mu)}
\xi^\nu .
\end{equation}
The residual bigram rule splits prediction into an ambiguous mean part and a
context-gated branch-specific residual:
\begin{equation}
h_i^{\mathrm{res}}(z)
=
\sum_\mu
\bar{\xi}_i^{\mu+1}
m_\mu(z)|m_\mu(z)|^2
+
\sum_\mu
\left(
\xi_i^{\mu+1}
-
\bar{\xi}_i^{\mu+1}
\right)
m_\mu(z)|m_{\mu-1}(z)|^2 .
\end{equation}

\paragraph{Ridge-fitted quartic control.}
As a fitted upper-bound control, we also use all features
$m_a|m_b|^2$ and fit their coefficients by ridge regression:
\begin{equation}
h_i^{\mathrm{ridge}}(z)
=
\sum_\nu
\xi_i^\nu
\sum_{a,b}
W_{\nu ab}
m_a(z)
|m_b(z)|^2 ,
\end{equation}
where $W_{\nu ab}$ is fit from blended anchor states near stored transitions.

\section{Empirical Success-Rate Classification}

\paragraph{Goal.}
We want to classify which static fourth-order Kuramoto learning rules actually
retrieve crossing sequences autonomously. The key question is not only whether
a rule stores the intended transitions, but whether the resulting trajectory
visits the correct patterns in the correct temporal order and chooses the
correct branch at shared states.

\paragraph{Protocol.}
For each benchmark topology and each random pattern realization, we initialize
the network as a small perturbation of the first pattern of one target sequence,
\begin{equation}
    z(0) \approx \xi^{(s,1)} ,
\end{equation}
and then evolve the autonomous Kuramoto dynamics with no external drive.  Each
stored branch is tested as a separate retrieval trial.

\paragraph{Rules to compare.}
The empirical comparison should include:
\begin{itemize}
    \item current-only baselines: self, forward $(\mu\to\mu+1)$, skip-two
    $(\mu\to\mu+2)$, and their mixtures;
    \item context rules: bigram, multilag, coherent bigram, coherent trigram,
    and residual bigram;
    \item fitted or oracle controls: ridge-fitted quartic rule and, optionally,
    a lifted Markov-1 rule with branch-specific copies of shared states.
\end{itemize}

\paragraph{Primary success metric.}
The main success criterion is peak-order retrieval: for the target sequence
$(\alpha_1,\ldots,\alpha_L)$, the maximum-overlap times must satisfy
\begin{equation}
    t^\ast_{\alpha_1}
    <
    t^\ast_{\alpha_2}
    <
    \cdots
    <
    t^\ast_{\alpha_L},
    \qquad
    t^\ast_{\alpha}
    =
    \argmax_t |m_\alpha(z(t))| .
\end{equation}
For crossing benchmarks we additionally record whether the selected outgoing
branch is the intended one.  Thus the most important reported quantity is the
branch-correct peak-order success rate.

\paragraph{Success-rate estimate.}
For each rule and benchmark we report
\begin{equation}
\widehat{p}_{\mathrm{succ}}
=
\frac{1}{N_{\mathrm{trials}}}
\sum_{r=1}^{N_{\mathrm{trials}}}
\mathbf{1}\!\left\{
\text{trial } r \text{ satisfies the retrieval criterion}
\right\}.
\end{equation}
Confidence intervals should be shown across random pattern realizations, e.g. by
Wilson intervals for binomial success rates.

\paragraph{Experiments to finish the report.}
The empirical section should show:
\begin{itemize}
    \item a success-rate table or heatmap across the first three benchmarks;
    \item example trajectories showing Markov-1 failure and context-rule success;
    \item a branch-choice diagnostic at the shared state, comparing correct and
    wrong successor drive;
    \item robustness sweeps over $N$, corruption $\rho$, gain $g$, initial jitter,
    and noise;
    \item a memory-load experiment where unrelated distractor sequences are added;
    \item a short comparison showing that ridge is an upper-bound control, while
    the coherent trigram is the simplest interpretable rule that still solves the
    shared-segment benchmark.
\end{itemize}


\section{Related Work}

\paragraph{Kuramoto and oscillatory associative-memory background.}
Oscillatory associative-memory models provide the dynamical foundation for representing memories through phases rather than binary states.

\begin{itemize}
\item Aoyagi~\cite{aoyagi1995}: early phase-oscillator associative-memory model; establishes retrieval of phase-coded memories in recurrent oscillator networks.

\item Nishikawa, Hoppensteadt, and Lai~\cite{nishikawa2004}: analyzes oscillatory associative memory and shows how additional dynamical structure can stabilize retrieval states.
\item Zhao, Li, and Xue~\cite{zhao2020}: studies stability of Hebbian Kuramoto associative-memory networks with second-order coupling; provides direct theoretical support for Kuramoto-style associative retrieval.
\item Fu, Li, Lin, and Zhao~\cite{fu2025}: introduces higher-order self-dynamics in Kuramoto-type associative memory and shows enlarged retrieval basins while preserving memory capacity; focuses primarily on static memory retrieval rather than sequence disambiguation.

\end{itemize}

\paragraph{Dense associative-memory background.}
Dense associative memory extends classical Hopfield retrieval by replacing pairwise interactions with stronger nonlinear or higher-order interactions, thereby improving pattern separation and storage capacity.

\begin{itemize}
\item Krotov and Hopfield~\cite{krotov2016dense}: introduces dense associative memory and demonstrates that higher-order nonlinear interactions can substantially increase storage capacity and suppress crosstalk.

\item Demircigil et al.~\cite{demircigil2017}: establishes exponential-capacity associative memory for sufficiently strong nonlinear interaction functions, providing theoretical motivation for high-order overlap terms.
\item Chaudhry et al.~\cite{chaudhry2024long}: extends dense associative-memory ideas to sequence storage; nonlinear overlap separation improves sequence capacity and pseudoinverse-style learning supports correlated stored patterns.

\end{itemize}

\paragraph{Sequential-memory background.}
Sequential associative-memory models introduce directed temporal structure so that retrieval evolves from one stored pattern to the next rather than converging only to static attractors.

\begin{itemize}
\item Sompolinsky and Kanter~\cite{sompolinsky1986temporal}: classical asymmetric Hopfield-style sequence memory; stores directed temporal associations through asymmetric synaptic weights.

\item Kleinfeld~\cite{kleinfeld1986sequential}: separates stabilizing associative connections from delayed transition connections; provides an early mechanism for generating ordered sequences of attractor states.
\item Kawamura and Okada~\cite{kawamura2002transient}: derives macroscopic transient dynamics for sequential associative-memory networks and supports an overlap-based description of sequence retrieval.
\item Mimura, Kawamura, and Okada~\cite{mimura2004path}: applies path-integral analysis to associative memories storing finite limit-cycle sequences and studies retrieval dynamics and storage capacity.
\item Gillett, Pereira, and Brunel~\cite{gillett2020tah}: analyzes sequential activity in rate and spiking networks learned through temporally asymmetric Hebbian plasticity, including analytic characterization of sequence capacity.
\item Chaudhry et al.~\cite{chaudhry2024long}: provides the closest modern Hopfield-style sequence-memory formulation to the present setting, combining nonlinear associative interactions with explicitly directed sequence transitions.

\end{itemize}

\paragraph{Context, crossings, and sequence-disambiguation background.}
A current-state-only sequence rule becomes ambiguous when two stored trajectories contain the same present pattern but require different successors.  A separate line of work therefore introduces contextual state, traces, bias signals, or higher-order temporal representations to resolve such crossings.

\begin{itemize}
\item Levy and Wu~\cite{levy1996context}: analyzes local context codes for sequence prediction and shows how the lifetime of contextual information constrains sequence-memory capacity.

\item Sohal and Hasselmo~\cite{sohal1998}: studies sequence disambiguation in a hippocampal CA3 model and shows that temporally structured modulation can distinguish sequences containing shared intermediate states.
\item Kawamura, Okada, and Hirai~\cite{kawamura1999selective}: studies one-to-many associative recall and shows that selecting the correct successor requires additional seed or contextual information.
\item Katahira et al.~\cite{katahira2007branching}: directly studies branching sequence retrieval; a shared state can be driven toward the intended branch by additional common and bias inputs.
\item Samura, Hattori, and Ishizaki~\cite{samura2008}: studies linked sequences containing shared events and shows that cooperation between autoassociative and heteroassociative mechanisms can support sequence disambiguation.
\item Cui, Ahmad, and Hawkins~\cite{cui2016}: develops a high-order temporal sequence-memory model in which identical feedforward inputs can correspond to distinct internal states depending on temporal context.
\item Hawkins and Ahmad~\cite{hawkins2016synapses}: proposes a biological sequence-memory mechanism based on sparse activity and dendritic context; the same feedforward pattern can participate in different sequence states depending on recent history.
\item Martinez Mayorquin et al.~\cite{martinez2019}: uses low-pass-filtered synaptic traces to disambiguate sequences containing shared elements; this is particularly close to augmenting the current state with a delayed context variable.

\end{itemize}

Taken together, these literatures motivate the rule studied here.  Kuramoto associative memory provides the phase-based retrieval dynamics; dense associative memory motivates nonlinear overlap interactions; asymmetric sequence-memory models provide directed transitions between stored patterns; and context-based sequence-memory models show why an additional history-dependent variable is required when stored trajectories cross.  Our simplest interpretable rule that succeeds across the first three benchmarks is therefore a lag-context, coherent-trigram rule: the next recalled pattern is driven by the current overlap, but disambiguated by recent overlap history.

\printbibliography
\end{document}
